% 3.4 章 · 边值问题与谱方法
% 这一章的两条主线：
%   1) 打靶法把边值问题变成初值问题，代价是一个叫"灵敏度"的放大因子；
%   2) 谱方法用 N+1 个 Chebyshev 节点换指数收敛，代价是一个稠密、条件数按 N^2 长的微分矩阵。
% 两条主线的共同点是：**失效的位置由舍入决定，不由截断决定**。
% 断言附表由 tools/build_book.py 从实际运行结果生成，\input 进来，不手抄。

\numlabchapter{边值问题与谱方法}{3.4}

\section{这一章要解决什么}

两点边值问题
\begin{equation}
  u''+p(x)u'+q(x)u=r(x),\qquad x\in[a,b],
  \label{eq:bvp}
\end{equation}
配上一端一个边界条件（本章用 Dirichlet $u=g$ 与 Neumann $u'=g$）。它和初值问题
有一个本质差别：\textbf{条件给在两端，而积分只能从一端走}。这一个差别生出两条完全不同的路线：

\begin{enumerate}
  \item \textbf{打靶法}：猜一个左端斜率 $s=u'(a)$，把问题当初始问题积分到 $b$，
        解标量方程 $F(s)=u(b;s)-\beta=0$。它把边值问题压成一个根，
        代价是"猜错的 $s$ 会在右端被放大"。
  \item \textbf{离散化}（有限差分 / 谱配置 / tau）：把 $u$ 换成 $N+1$ 个未知数，
        连同边界条件一起解一个线性（或线性化后的）方程组。它没有放大问题，
        代价是矩阵的条件数与规模。
\end{enumerate}

本章要用实测数字把三件事说清楚：打靶法\textbf{什么时候不能用了}、有限差分的阶
\textbf{由谁决定}、以及谱方法的"指数收敛"\textbf{在哪里停下}。
记号：$u=\tfrac12\varepsilon_{\text{mach}}=2^{-53}\approx1.11\times10^{-16}$；
$L=b-a$；所有"实测"数字都来自 \texttt{tests/3.4/} 的断言报告。

\section{打靶法：灵敏度是唯一的判据，但不是充分的判据}

\begin{definition}[打靶灵敏度]
设 $u(x;s)$ 是初值问题 $u(a)=\alpha$、$u'(a)=s$ 的解。定义
\begin{equation}
  S:=\left|\frac{\partial u(b;s)}{\partial s}\right|.
  \label{eq:sens}
\end{equation}
$S$ 就是"左端斜率的误差传到右端时被放大多少倍"。
\end{definition}

\begin{theorem}[变分方程给出 $S$，不必做差分]
\label{thm:variational}
令 $y=(u,u')^{\mathsf T}$，$J(x)=\partial f/\partial y$。则 $z=\partial y/\partial s$ 满足
\begin{equation}
  z'=J(x,u(x))z,\qquad z(a)=(0,1)^{\mathsf T},
  \label{eq:variational}
\end{equation}
且 $\partial u(b)/\partial s=z_1(b)$。
\end{theorem}
\begin{proof}
对 $s$ 求导并与 $y'=f(x,y)$ 交换次序即得 \eqref{eq:variational}。
\end{proof}

\begin{remark}[实测：$u''=k^2u$、$u(a)=0$ 的闭式就是 $\sinh(kL)/k$]
\label{rem:sens-closed}
该方程的通解是 $Ae^{kx}+Be^{-kx}$，$u(a)=0$ 给出 $A=-Be^{-2ka}$，
于是 $u(b;s)=\frac{s}{k}\sinh(kL)$，即 $S=\sinh(kL)/k$。
本章用变分方程积分出的 $S$ 与之对拍：$kL\le20$ 时相对偏差 $\le1.6\times10^{-9}$，
$kL=40$ 时 $5.3\times10^{-8}$（$S$ 自己已经 $2.9\times10^{15}$，变分方程在积一个 $e^{40}$ 的增长模态）。
\end{remark}

\begin{theorem}[线性问题 Newton 一步，非线性问题二次收敛]
线性问题里 $F(s)=u(b;s)-\beta$ 对 $s$ 是\textbf{仿射}的，故 Newton 一次到；
非线性问题里 $F''\ne0$，残差按 $|F_{k+1}|\le C|F_k|^2$ 下降。
\end{theorem}
\begin{proof}
线性：$u(b;s)=u_0+s\,z_1(b)$ 是 $s$ 的一次函数。非线性：标准 Newton 二次收敛。
\end{proof}

\begin{remark}[实测：仿射一步、Bratu 二次]
\texttt{smooth}、\texttt{neumann}、\texttt{layer} 三个线性算例都在 2 次函数求值内收敛
（1 步 Newton + 1 次确认）。Bratu（$u''=1.5u^2$、$u=4/(1+x)^2$）的残差史是
$2.7\times10^{1}\to8.5\to1.9\to1.8\times10^{-1}\to2.1\times10^{-3}\to3.0\times10^{-7}\to7.3\times10^{-13}$，
7 步到 $10^{-12}$，误差 $2.4\times10^{-11}$。
\end{remark}

\begin{theorem}[衰减模态的失效是舍入，不是截断]
\label{thm:decay-fail}
考虑 $u''=k^2u$、$u(0)=1$、$u(L)=0$，精确解 $u=\sinh(k(L-x))/\sinh(kL)$。
它由两个指数模态 $e^{\pm kx}$ 相减得到，而 $e^{kL}$ 远大于解本身，
于是相对误差被放大 $e^{kL}$ 倍：
\begin{equation}
  \text{err}\sim C\,u\,\varepsilon_{\text{mach}}e^{kL}.
  \label{eq:decay-err}
\end{equation}
这个放大与步长 $h$ 无关——它不是离散误差。
\end{theorem}
\begin{proof}[证明思路]
把初始斜率写成 $s=s^*+\delta$。$u(b;s)$ 对 $s$ 的灵敏度是 $S=\sinh(kL)/k$
（定理 \ref{thm:variational} 与注 \ref{rem:sens-closed}），
而 $u(b;s^*)=0$，所以右端的偏差是 $S\delta$；为把它压到 $u$ 的量级，
$\delta$ 必须小到 $u/S$，这已经低于 $s^*$ 本身的可分辨精度 $u\,|s^*|$
（因为 $|s^*|=k\coth(kL)\to k$ 而 $S/k=e^{kL}/(2kL)\cdot2kL$ 量级）。
两者之比 $e^{kL}$ 就是丢掉的有效位数。
\end{proof}

\begin{remark}[实测：同一个 $S$，两种命运]
$k=20$、$L=2$（$kL=40$）时，增长算例 \texttt{layer}（$u(0)=0,u(L)=1$）与
衰减算例 \texttt{decay} 的实测灵敏度\textbf{完全相同}（$5.885\times10^{15}$），
误差却是 $4.8\times10^{-10}$ 与 $1.7\times10^{1}$，差 10 个量级。
把精确斜率喂进去、步长从 $h=8\times10^{-3}$ 缩到 $2.5\times10^{-4}$，
衰减算例的误差仍在 $0.6\sim26$ 之间不下降：\eqref{eq:decay-err} 是舍入放大，不是截断。
\end{remark}

\begin{remark}[实测：失效长度的二分阈值]
取判据"打靶误差 $>0.5$"，对 $k=20$ 二分实测 $L^*=1.8750$，
而舍入估计 $\ln(1/u)/k=1.8373$（用 $u=\varepsilon_{\text{mach}}/2$），
相对差 $2\%$。误差随 $kL$ 的增长也实测为指数律：半对数斜率 $0.89$（理论 1），
$R^2=0.955$。
\end{remark}

\section{有限差分：整体阶由边界格式决定}

\begin{theorem}[内部二阶中心差商 + 边界行]
在等距网格 $x_i=a+ih$ 上取
\begin{equation}
  \frac{u_{i+1}-2u_i+u_{i-1}}{h^2}+p_i\frac{u_{i+1}-u_{i-1}}{2h}+q_iu_i=r_i,
  \qquad i=1,\dots,n-1,
  \label{eq:fd}
\end{equation}
两端各补一行边界条件。截断误差是 $O(h^2)$。
\end{theorem}

\begin{theorem}[边界差商的阶会拖住整体阶]
\label{thm:bc-order}
若 Neumann 边界用一阶单侧差商 $\frac{u_n-u_{n-1}}{h}=g$，则\textbf{整体}误差是 $O(h)$，
即使内部是二阶；用二阶单侧差商则整体回到 $O(h^2)$。
误差不局限在边界附近：一阶边界下"掐掉两端 $25\%$ 再统计"得到的阶仍是 $1.00$。
\end{theorem}
\begin{proof}[证明思路]
把边界行的截断误差 $\tau_b=O(h)$ 看成右端的一个扰动。
差分算子 $A$ 的逆的范数是 $O(h^{-2})=O(n^2)$ 量级（实测 $\kappa_\infty\sim n^2$），
但扰动只作用在一行上，Green 函数的影响是 $O(h)$ 且沿整个区间铺开，
于是它进入解的量级是 $O(h)$，与内部截断同阶竞争。
\end{proof}

\begin{remark}[实测：0.9983 与 1.9756]
混合边界算例（$u=\sin(\pi x/2)$，$u(0)=0$、$u'(1)=0$）：
一阶边界格式在 $n=32..512$ 上整体阶 $0.9983$、内部阶 $0.9981$；
二阶边界格式在 $n=128..1024$ 上整体阶 $1.9756$、内部阶 $1.9780$。
\textbf{报阶必须写窗口}：同一批二阶数据在 $n=16..128$ 上只有 $1.86$。
\end{remark}

\begin{remark}[边界残差的两个口径]
一阶边界格式、$n=64$ 时，"格式自己那一行"的残差是 $3.3\times10^{-15}$，
而把数值解当真函数看、边界条件的真实违反量是 $1.93\times10^{-2}$——差 12 个量级。
前者恒为舍入级（解就是照那一行解出来的），后者才是精度的证据。
同理，线性系统的残差与解的误差无关：$n=32$ 时残差 $8.7\times10^{-14}$、误差 $1.57\times10^{-4}$，
比值 $1.8\times10^{9}$。
\end{remark}

\section{谱方法：Chebyshev 节点、微分矩阵与两条收敛路线}

\begin{definition}[Chebyshev--Lobatto 节点]
$x_j=\cos(j\pi/N)$，$j=0,\dots,N$（降序，含两端）。
它是 $\theta$ 均匀网格的余弦投影，密度 $\rho(x)=N/(\pi\sqrt{1-x^2})$。
\end{definition}

\begin{theorem}[配置微分矩阵]
\label{thm:dmatrix}
矩阵 $D_N$ 的非对角元为
\begin{equation}
  D_{ij}=\frac{c_i}{c_j}\frac{(-1)^{i+j}}{x_i-x_j},\qquad c_0=c_N=2,\ c_k=1,
  \label{eq:dmatrix}
\end{equation}
对角元取 $D_{ii}=-\sum_{j\ne i}D_{ij}$。于是 $D_N\mathbf{1}=0$（零空间是常数）、
$\|D_N\|_\infty=N^2$、$D_{00}=(2N^2+1)/6$。
\end{theorem}
\begin{proof}[证明思路]
\eqref{eq:dmatrix} 由 Lagrange 插值多项式的导数在节点上取值得到；
行和为零是因为常数函数的导数恒为 $0$；$\|D_N\|_\infty=N^2$ 是已知的精确等式
（本章逐位复核过：$N=4..64$ 全部吻合）。
\end{proof}

\begin{remark}[实测：$\|D_N\|_\infty=N^2$ 逐位吻合，$D_{00}=(2N^2+1)/6$ 到 $10^{-13}$]
$N=8$：$\|D_N\|_\infty=64.00$、$D_{00}=21.5$；$N=16$：$256.00$、$85.5$；
$N=64$：$4096.00$、$1365.17$。$\|D_N^2\|_\infty\approx N^4/3$（$N=32$ 实测 $3.49\times10^{5}$）。
\end{remark}

\begin{theorem}[指数收敛与代数收敛]
\label{thm:conv}
若 $u$ 在含 $[a,b]$ 的带形区域内解析，则 Chebyshev 插值（从而配置解）误差按
$e^{-\lambda N}$ 衰减；若 $u$ 只有 $C^s$ 光滑度，则按 $N^{-s}$ 衰减，
而\textbf{插值}误差是 $N^{-(s+1)}$：微分吃掉一阶。
\end{theorem}

\begin{remark}[实测：斜率、地板、以及撞上地板的 $N$]
解析算例 \texttt{smooth}（$u=e^x\sin\pi x$）在 $N=8..18$ 上拟合半对数斜率 $-2.27$
（每节点约 $0.99$ 位，$R^2=0.998$），地板 $7.55\times10^{-15}$，
$N\ge18$ 起不再下降，最小误差出现在 $N=24$，$N=64$ 时回到 $9.3\times10^{-14}$。
非解析算例：$|x|^3$ 的配置解率 $1.99$、插值率 $3.00$；$|x|^5$ 是 $4.02$ 与 $5.01$。
两者都只有幂律——判别指数/幂律不能只看 $R^2$（\texttt{tanh} $k=3$ 的窗口里两种模型都拟合得不错），
要看误差比 $\mathrm{err}(N)/\mathrm{err}(N+2)$：指数收敛的比值整体上抬（实测 $39\to76\to197$），
幂律收敛的比值压向常数（实测 $1.56\to1.14$）。
\end{remark}

\begin{remark}[实测：舍入地板是 $N^4u$]
把\textbf{精确解}插值后做二阶谱微分，残差在 $N=20/24/32$ 上是
$6.8\times10^{-12}/2.1\times10^{-11}/1.1\times10^{-10}$，
而同一批节点上的插值误差只有 $3.3\times10^{-15}$；$N^4u$ 给出
$5.5\times10^{-12}/1.2\times10^{-11}/3.6\times10^{-11}$。
配置法内点残差则恒在舍入级（$N=8$ 时 $1.7\times10^{-13}$），
而节点\textbf{之间}的连续残差是 $0.38$——比解误差 $1.8\times10^{-4}$ 还大 3 个量级。
\end{remark}

\begin{theorem}[tau 法与配置法的差别在边界行]
配置法在节点上逐点满足 ODE（$N-1$ 个内点），再加两行边界；
tau 法在系数空间满足前 $N-1$ 个方程，再补两行边界。
两者给出同一个多项式，但对 Neumann 边界的写法不同（本章的 tau 只实现两端 Dirichlet）。
\end{theorem}
\begin{proof}[证明思路]
配置方程 $D_N^2u+\operatorname{diag}(q)u=r$ 与系数空间方程在多项式空间上等价，
差别只在被丢弃的方程与边界行的组合方式。
\end{proof}

\begin{remark}[实测：两条路线互差 $9.7\times10^{-15}$]
$N=20$ 时配置解误差 $8.5\times10^{-15}$、tau 解 $6.9\times10^{-15}$，两者之差 $9.7\times10^{-15}$；
条件数同量级（$1.6\times10^{4}$ 与 $2.6\times10^{4}$），矩阵密度 $0.91$ 与 $0.37$。
\end{remark}

\begin{remark}[区间映射的伸缩因子]
$D_N$ 是参考区间 $[-1,1]$ 上的矩阵。问题在 $[a,b]$ 上时
$\mathrm{d}/\mathrm{d}x=s\,\mathrm{d}/\mathrm{d}\xi$、
$\mathrm{d}^2/\mathrm{d}x^2=s^2\,\mathrm{d}^2/\mathrm{d}\xi^2$，$s=2/(b-a)$。
$[0,1]$ 上忘记 $s$ 会解出偏差 $O(1)$ 的东西（实测 $>0.5$）。
识别"这是不是 Lobatto 网格"也必须允许平移缩放，否则会把 $[0,1]$ 上的谱网格
当成普通网格，边界残差当场量出 $\pi/2$ 这种荒唐数字。
\end{remark}

\section{四个可视化该怎么看}

\begin{enumerate}
  \item \texttt{shooting\_\allowbreak{}scan}（招牌）：一族初始斜率的轨迹
        $u(x;s_k)$，$s_k=s_{\text{ref}}\bigl(1+\text{span}(2k/(m-1)-1)\bigr)$。
        时间轴是斜率序号 $j$。副轴是对数坐标下的打靶残差 $|F(s_k)|$。
        选 \texttt{decay} 时看的是"参考斜率附近的微小扰动在右端被放大成 $10^3$ 量级"，
        选 \texttt{layer} 时同一族轨迹规规矩矩——两次点击就能看见定理 \ref{thm:decay-fail}。
  \item \texttt{convergence\_\allowbreak{}rate}：有限差分与谱配置的误差并排。
        左轴半对数（指数收敛是直线）、右轴双对数（幂律是直线），
        两条\textbf{实测}拟合线叠加在数据上，地板线单独画，插值误差作为负对照。
        把 \texttt{bc\_\allowbreak{}order} 从 $2$ 切到 $1$，差分那条曲线整体抬高一个阶。
  \item \texttt{node\_\allowbreak{}clustering}：节点位置、间距、解析密度，
        以及"节点上插值误差恒为 $0$、节点之间才是误差"的那条曲线。
        参数 $N$ 拖大，看间距比按 $0.6366N$ 长、配置解误差与插值误差的差距保持一阶。
  \item \texttt{diff\_\allowbreak{}matrix\_\allowbreak{}structure}：$\log_{10}|D_N|$ 热力图、
        行和残差（应恒为 $0$，零空间是常数）与奇异值谱。
        注意 $\sigma_{\min}\sim10^{-17}$ 是零空间噪声，有意义的是挖掉零空间后的
        $\kappa^+$（$N=8..64$ 实测 $35.1\to1988$，指数 $1.95$）；
        而配置矩阵 $\kappa_\infty(L)$ 从 $5.5\times10^{2}$（$N=8$）长到 $2.8\times10^{9}$（$N=256$）。
\end{enumerate}

\section{本章的算法路线（P2/P3 合规说明）}

\begin{itemize}
  \item 全部线性系统走自研 \texttt{numcore.linalg.lu} 的
        \texttt{lu\_decompose}/\texttt{solve\_lu}（部分主元），
        奇异值走自研 Jacobi SVD；\texttt{numpy.linalg} 一次都没用；
  \item 拟合不用 \texttt{numpy.polyfit}：半对数/双对数的最小二乘直线是自写的
        \texttt{\_line\_fit} 与 $R^2$，窗口写进每条断言的 \texttt{theory}/\texttt{detail}；
  \item 打靶的变分方程与主方程共用 \texttt{numcore.ode.rk} 的 \texttt{rk\_step}，
        不做数值差分；Newton 的收敛判据用\textbf{相对步长}；
  \item "失效长度"这类阈值一律二分实测（\texttt{find\_threshold}），
        理论值只用来对账；
  \item 每一条"指数收敛"的结论都带\textbf{实测斜率 + 舍入地板 + 撞上地板的 $N$}，
        每一条"阶"的结论都带\textbf{拟合窗口}。
\end{itemize}

% 本文件由 tools/build_book.py 从 dist/reports/assertions.json 生成，请勿手改。
\begin{longtable}{@{}p{0.06\textwidth}p{0.32\textwidth}p{0.22\textwidth}p{0.16\textwidth}@{}}
\caption{章节 3.4 数值断言（生成时间 2026-10-19T16:08:50-0400）}\\
\toprule
类别 & 断言 & 理论值 & 实测值 \\
\midrule
\endfirsthead
\toprule
类别 & 断言 & 理论值 & 实测值 \\
\midrule
\endhead
\bottomrule
\endlastfoot

A & 帧契约：角点 $\times$ 时间端点的合法帧数 & 所有参数角点与时间端点都必须建出有限、\allowbreak{}合规的帧 & 4 个 viz，共 60 帧、\allowbreak{}711 个数组段，全部有限且 meta 全是有限标量 \checkmark \\
B & 帧内灵敏度 vs 直接重算(相对) & $\le$ 1.00000e-12 & 0 \checkmark \\
B & 灵敏度 vs 闭式 sinh(kL)/k(相对) & $\le$ 1.00000e-06 & $2.013\times 10^{-9}$ \checkmark \\
B & 帧内参考斜率 vs 直接重算(相对) & $\le$ 1.00000e-12 & 0 \checkmark \\
A & viz1 的数字对拍(decay, k=10, L=1) & S=\allowbreak{}sinh(\allowbreak{}kL)/\allowbreak{}k=\allowbreak{}1101.32；s\_ref=\allowbreak{}-k*coth(\allowbreak{}kL) & 帧内 S=\allowbreak{}1.101323e+03，直接重算 1.101323e+03(\allowbreak{}相对 0.0e+00)；闭式 1.101323e+03(\allowbreak{}相对 2.0e-09)；s\_ref=\allowbreak{}-10.000000 \checkmark \\
B & 帧内地板 vs 直接重算(相对) & $\le$ 1.00000e-12 & 0 \checkmark \\
B & 帧内半对数斜率/R\textasciicircum{}2 vs 直接重算 & $\le$ 1.00000e-09 & 0 \checkmark \\
A & viz2 的数字对拍(smooth) & 半对数斜率由未触地板的粗端窗口拟合；地板=\allowbreak{}曲线最小误差 & 斜率 -2.0894(\allowbreak{}重算 -2.0894)，R\textasciicircum{}2=\allowbreak{}0.99491，地板 7.550e-15(\allowbreak{}重算 7.550e-15)，窗口 N=\allowbreak{}4..16(\allowbreak{}head=\allowbreak{}7) \checkmark \\
B & 聚集比 vs 直接重算 & $\le$ 1.00000e-12 & 0 \checkmark \\
B & 密度比 vs 直接重算 & $\le$ 1.00000e-12 & 0 \checkmark \\
B & 密度比 vs 解析 1/sqrt(0.19) & $\le$ 1.00000e-12 & 0 \checkmark \\
A & viz3 的数字对拍(abs3, N=24) & 聚集比 0.6366N；密度比 1/\allowbreak{}sqrt(\allowbreak{}0.19)=\allowbreak{}2.294；配置解比插值差一阶 & 聚集比 15.2571(\allowbreak{}重算 15.2571)；密度比 2.294157(\allowbreak{}解析 2.294157)；插值峰值 7.76e-05，配置解误差 8.39e-03 \checkmark \\
B & ||D||\_\allowbreak{}inf vs N\textasciicircum{}2(相对) & $\le$ 1.00000e-12 & $6.661\times 10^{-16}$ \checkmark \\
B & 行和残差 vs 直接重算 & $\le$ 1.00000e-18 & 0 \checkmark \\
B & kappa+ vs 直接重算 & $\le$ 1.00000e-12 & 0 \checkmark \\
B & D\_\allowbreak{}00 vs (2N\textasciicircum{}2+1)/6(相对) & $\le$ 1.00000e-13 & $1.163\times 10^{-15}$ \checkmark \\
A & viz4 的数字对拍(N=16, D) & ||D||\_inf=\allowbreak{}N\textasciicircum{}2；D$\cdot$1=\allowbreak{}0；kappa+ \textasciitilde{} N\textasciicircum{}2；朴素 kappa 无意义 & ||D||=\allowbreak{}256.000000(\allowbreak{}N\textasciicircum{}2=\allowbreak{}256)，行和残差 4.44e-15，朴素 kappa=\allowbreak{}1.36e+18，kappa+=\allowbreak{}129.956 \checkmark \\
B & 灵敏度 vs sinh(kL)/k(相对，kL$\le$20) & $\le$ 1.00000e-08 & $1.653\times 10^{-9}$ \checkmark \\
B & kL=40 时灵敏度自身的相对偏差 & $\le$ 1.00000e-06 & $5.245\times 10^{-8}$ \checkmark \\
A & 变分方程灵敏度 vs 闭式 sinh(kL)/k & S =\allowbreak{} sinh(\allowbreak{}kL)/\allowbreak{}k & kL$\le$20 最大相对偏差 1.653e-09；kL=\allowbreak{}40 时 S=\allowbreak{}5.884631e+15(\allowbreak{}闭式 5.884632e+15，相对 5.2e-08) \checkmark \\
B & 线性问题 Newton 收敛后的斜率(与另一次打靶互差，相对) & $\le$ 1.00000e-08 & $5.764\times 10^{-16}$ \checkmark \\
A & 线性问题 Newton 迭代次数 & 1 步(\allowbreak{}仿射 F)+ 1 次确认求值 & 最多 2 次求值；smooth: s=\allowbreak{}-1.155727e+00，2 次求值，err=\allowbreak{}1.07e-10；neumann: s=\allowbreak{}1.570796e+00，2 次求值，err=\allowbreak{}2.47e-13；layer: s=\allowbreak{}6.738253e-02，2 次求值，err=\allowbreak{}7.39e-11 \checkmark \\
A & Bratu 打靶 Newton 的二次收敛 & |F\_\{k+1\}| $\le$ 0.2$\cdot$|F\_k|\textasciicircum{}2(\allowbreak{}牛顿法的标准二次界，只查 |F|$\ge$1e-4 的段) & 残差史 2.69e+01 -> 8.52e+00 -> 1.92e+00 -> 1.82e-01 -> 2.12e-03 -> 3.03e-07 -> 7.25e-13 \checkmark \\
B & 两条路线的灵敏度一致(相对) & $\le$ 1.00000e-09 & 0 \checkmark \\
A & 增长模态 vs 衰减模态(同一 S) & S=\allowbreak{}sinh(\allowbreak{}kL)/\allowbreak{}k 相同，但衰减模态要在指数小的分量上做相减 & S=\allowbreak{}5.8846e+15；增长模态 err=\allowbreak{}4.82e-10，衰减模态 err=\allowbreak{}1.73e+01，相差 3.59e+10 倍 \checkmark \\
A & 精确斜率 + 加密步长下的衰减模态误差 & 舍入被 e\textasciicircum{}\{kL\} 放大：误差 \textasciitilde{} eps$\cdot$e\textasciicircum{}\{kL\}，与 h 无关 & ns=\allowbreak{}250..8000(\allowbreak{}h: 8.0e-3 $\to$ 2.5e-4)误差 6.02e-01\textasciitilde{}2.60e+01，kL=\allowbreak{}40，eps$\cdot$exp(\allowbreak{}kL)=\allowbreak{}2.6e+01 \checkmark \\
A & Dirichlet 两端：整体阶 & 2 & 2.00078 \checkmark \\
A & Dirichlet 两端：内部阶 & 2 & 2.00078 \checkmark \\
B & 离散边界残差(Dirichlet 被钉死) & $\le$ 1.00000e-11 & $2.316\times 10^{-12}$ \checkmark \\
A & Dirichlet：整体阶 = 内部阶 = 2 & 二阶中心差商 + 精确 Dirichlet 行 & 整体 2.0008，内部 2.0008；误差 4.331e-03 -> 1.688e-05 \checkmark \\
A & 混合边界(边界格式 1 阶)：整体阶 & 1 & 0.99946 \checkmark \\
A & 混合边界(边界格式 2 阶)：整体阶 & 2 & 1.97556 \checkmark \\
A & 混合边界：整体阶与内部阶 & 整体阶 =\allowbreak{} min(\allowbreak{}内部阶,\allowbreak{} 边界格式阶) & 一阶边界：整体 0.9995 /\allowbreak{} 内部 0.9994；二阶边界：整体 1.9756 /\allowbreak{} 内部 1.9780 \checkmark \\
B & 离散边界残差(格式自己那一行) & $\le$ 1.00000e-12 & $7.886\times 10^{-14}$ \checkmark \\
A & 边界残差两个口径的差距(一阶边界格式) & 离散行残差 \textasciitilde{} eps；真实违反量 \textasciitilde{} h$\cdot$|u''| & n=\allowbreak{}64: 离散 3.3e-15 /\allowbreak{} 真实 1.9e-02 /\allowbreak{} 误差 1.9e-02；n=\allowbreak{}128: 离散 1.8e-14 /\allowbreak{} 真实 9.6e-03 /\allowbreak{} 误差 9.6e-03；n=\allowbreak{}256: 离散 7.9e-14 /\allowbreak{} 真实 4.8e-03 /\allowbreak{} 误差 4.8e-03；两者相差 6.1e+10 倍 \checkmark \\
B & 线性系统残差(LU 回代) & $\le$ 1.00000e-09 & $2.762\times 10^{-11}$ \checkmark \\
A & 二阶边界格式下真实边界残差的衰减阶 & 3 & 3.00035 \checkmark \\
A & 残差与误差的比(有限差分) & 残差 \textasciitilde{} eps$\cdot$||A||，误差由截断决定，两者量级无关 & 最大线性残差 2.76e-11，误差 4.45e-05(\allowbreak{}n=\allowbreak{}64)-> 7.73e-07(\allowbreak{}n=\allowbreak{}512)；最小比值 2.8e+04 \checkmark \\
B & 每行非零元个数(三对角) & $\le$ 3.05 & 2.9922 \checkmark \\
A & FD 矩阵的稀疏结构与消元代价 & 三对角：每行 3 个非零；稠密 LU \textasciitilde{} n\textasciicircum{}3/\allowbreak{}3，追赶法 \textasciitilde{} 5n & 每行非零 ['2.88',\allowbreak{} '2.97',\allowbreak{} '2.99']；n=\allowbreak{}512 时稠密/\allowbreak{}带状 =\allowbreak{} 17476.3(\allowbreak{}n=\allowbreak{}32: 68.3) \checkmark \\
B & 打靶 vs 闭式解 & $\le$ 1.00000e-09 & $1.529\times 10^{-12}$ \checkmark \\
B & 有限差分(n=320)vs 闭式解 & $\le$ 1.00000e-04 & $4.671\times 10^{-6}$ \checkmark \\
B & 谱配置(N=32)vs 闭式解 & $\le$ 1.00000e-11 & $1.110\times 10^{-15}$ \checkmark \\
A & Bratu：三条路线的误差与 Newton 迭代数 & u=\allowbreak{}4/\allowbreak{}(\allowbreak{}1+x)\textasciicircum{}2 & 打靶 1.53e-12(\allowbreak{}7 次 F 求值)；FD n=\allowbreak{}320 4.67e-06(\allowbreak{}5 步)；配置 N=\allowbreak{}32 1.11e-15(\allowbreak{}5 步)；互差界 4.67e-06 \checkmark \\
A & 非线性 FD 的 Newton 残差地板 & 线性系统残差地板 \textasciitilde{} eps$\cdot$||A|| \textasciitilde{} eps/\allowbreak{}h\textasciicircum{}2 & n=\allowbreak{}80: 残差 9.63e-12(\allowbreak{}eps/\allowbreak{}h\textasciicircum{}2=\allowbreak{}7.1e-13)；n=\allowbreak{}160: 残差 3.72e-11(\allowbreak{}eps/\allowbreak{}h\textasciicircum{}2=\allowbreak{}2.8e-12)；n=\allowbreak{}320: 残差 1.94e-10(\allowbreak{}eps/\allowbreak{}h\textasciicircum{}2=\allowbreak{}1.1e-11) \checkmark \\
C & 衰减模态打靶失效长度 L* & 1.83684 & 1.87488 \checkmark \\
C & 失效长度 L* 的两侧 & L* =\allowbreak{} ln(\allowbreak{}1/\allowbreak{}u)/\allowbreak{}k =\allowbreak{} 1.8368 & 二分 1.8749(\allowbreak{}区间 [1.8748,\allowbreak{} 1.8750]，13 次迭代，15 次谓词调用)；kL* =\allowbreak{} 37.50 \checkmark \\
C & 衰减模态打靶误差的指数律 & err \textasciitilde{} C$\cdot$e\textasciicircum{}\{kL\}(\allowbreak{}舍入被灵敏度放大)，半对数斜率 =\allowbreak{} 1 & kL=\allowbreak{}[5.0,\allowbreak{} 10.0,\allowbreak{} 15.0,\allowbreak{} 20.0,\allowbreak{} 25.0,\allowbreak{} 30.0] 对应 err=\allowbreak{}['1.19e-13',\allowbreak{} '1.92e-12',\allowbreak{} '3.54e-11',\allowbreak{} '8.21e-10',\allowbreak{} '2.34e-06',\allowbreak{} '5.24e-04']；拟合斜率 0.8927，R\textasciicircum{}2=\allowbreak{}0.95547 \checkmark \\
C & 同一格式两个窗口的实测阶 & 二阶格式的渐近阶 2；前渐近区会偏低 & n=\allowbreak{}16..128 得 1.7724；n=\allowbreak{}256..1024 得 1.9842 \checkmark \\
C & FD 矩阵条件数 \textasciitilde{} n\textasciicircum{}2 & cond\_inf \textasciitilde{} n\textasciicircum{}2(\allowbreak{}差分算子 \textasciitilde{} (\allowbreak{}1/\allowbreak{}h\textasciicircum{}2)$\cdot$三对角) & n=\allowbreak{}[32.0,\allowbreak{} 64.0,\allowbreak{} 128.0,\allowbreak{} 256.0] 对应 cond=\allowbreak{}['1.025e+03',\allowbreak{} '4.097e+03',\allowbreak{} '1.639e+04',\allowbreak{} '6.554e+04']；拟合指数 1.9996(\allowbreak{}R\textasciicircum{}2=\allowbreak{}1.00000) \checkmark \\
B & 变换-求值往返(值尺度 \textasciitilde{}8) & $\le$ 7.10543e-15 & $6.828\times 10^{-15}$ \checkmark \\
B & x\_\allowbreak{}j = cos(j pi/N)(精确到舍入) & $\le$ 1.00000e-15 & 0 \checkmark \\
A & Chebyshev 往返变换与节点定义 & 互逆；x\_j=\allowbreak{}cos(\allowbreak{}j pi/\allowbreak{}N)，j=\allowbreak{}0..N(\allowbreak{}降序、\allowbreak{}含端点) & 往返最大偏差 6.828e-15；节点最大偏差 0.000e+00(\allowbreak{}N=\allowbreak{}4..64) \checkmark \\
B & 一阶导数：系数空间 vs D\_\allowbreak{}N & $\le$ 1.00000e-10 & $1.883\times 10^{-13}$ \checkmark \\
B & 二阶导数：系数空间 vs D\_\allowbreak{}N\textasciicircum{}2 & $\le$ 1.00000e-08 & $5.877\times 10^{-11}$ \checkmark \\
A & 系数空间微分 vs 微分矩阵 & 系数空间求导 M1 与 D\_N 一致、\allowbreak{}M2 与 D\_N\textasciicircum{}2 一致(\allowbreak{}含 c\_0=\allowbreak{}c\_N=\allowbreak{}2 的因子) & N=\allowbreak{}24：一阶偏差 1.88e-13，二阶偏差 5.88e-11；再作用一次的自洽性偏差 2.55e+01 \checkmark \\
B & 对 m$\le$N 的 x\textasciicircum{}m 积分误差 & $\le$ 1.00000e-13 & $4.441\times 10^{-16}$ \checkmark \\
A & Clenshaw--Curtis 的代数精度 & $\int$\_\{-1\}\textasciicircum{}\{1\} x\textasciicircum{}m dx =\allowbreak{} 2/\allowbreak{}(\allowbreak{}m+1)(\allowbreak{}m 偶)，否则 0；代数精度 =\allowbreak{} N(\allowbreak{}N+2 处失效) & m$\le$N 的最大偏差 4.441e-16；N=\allowbreak{}4: m=\allowbreak{}N+2 处偏差 1.9e-02，sum(\allowbreak{}w)=\allowbreak{}2.00000000000000；N=\allowbreak{}8: m=\allowbreak{}N+2 处偏差 7.2e-05，sum(\allowbreak{}w)=\allowbreak{}2.00000000000000；N=\allowbreak{}16: m=\allowbreak{}N+2 处偏差 3.1e-08，sum(\allowbreak{}w)=\allowbreak{}2.00000000000000；N=\allowbreak{}32: m=\allowbreak{}N+2 处偏差 5.7e-14，sum(\allowbreak{}w)=\allowbreak{}2.00000000000000 \checkmark \\
B & ||D\_\allowbreak{}N||\_\allowbreak{}inf / N\textasciicircum{}2 - 1 & $\le$ 1.00000e-12 & $8.438\times 10^{-15}$ \checkmark \\
B & 行和(D\_\allowbreak{}N$\cdot$1 = 0) & $\le$ 1.00000e-10 & $2.274\times 10^{-13}$ \checkmark \\
B & D\_\allowbreak{}00 vs (2N\textasciicircum{}2+1)/6(相对) & $\le$ 1.00000e-13 & $1.082\times 10^{-14}$ \checkmark \\
B & ||D\_\allowbreak{}N\textasciicircum{}2||\_\allowbreak{}inf / (N\textasciicircum{}4/3) - 1 & $\le$ 0.02 & 0.015625 \checkmark \\
A & 微分矩阵的结构体检 & D 稠密(\allowbreak{}density=\allowbreak{}1)、\allowbreak{}D$\cdot$1=\allowbreak{}0、\allowbreak{}D\_00=\allowbreak{}(\allowbreak{}2N\textasciicircum{}2+1)/\allowbreak{}6、\allowbreak{}||D||\_inf=\allowbreak{}N\textasciicircum{}2、\allowbreak{}||D\textasciicircum{}2||\_inf$\approx$N\textasciicircum{}4/\allowbreak{}3 & N=\allowbreak{}8: ||D||=\allowbreak{}64(\allowbreak{}N\textasciicircum{}2=\allowbreak{}64)，行和 5.0e-16，||D\textasciicircum{}2||=\allowbreak{}1.344e+03；N=\allowbreak{}16: ||D||=\allowbreak{}256(\allowbreak{}N\textasciicircum{}2=\allowbreak{}256)，行和 4.4e-15，||D\textasciicircum{}2||=\allowbreak{}2.176e+04；N=\allowbreak{}32: ||D||=\allowbreak{}1024(\allowbreak{}N\textasciicircum{}2=\allowbreak{}1024)，行和 2.8e-14，||D\textasciicircum{}2||=\allowbreak{}3.492e+05；N=\allowbreak{}64: ||D||=\allowbreak{}4096(\allowbreak{}N\textasciicircum{}2=\allowbreak{}4096)，行和 2.3e-13，||D\textasciicircum{}2||=\allowbreak{}5.591e+06；最小密度 0.9877；D\_00 与 (\allowbreak{}2N\textasciicircum{}2+1)/\allowbreak{}6 最大相对偏差 1.08e-14 \checkmark \\
B & 配置法内点残差 & $\le$ 1.00000e-08 & $5.157\times 10^{-12}$ \checkmark \\
A & 配置法的残差三口径 & 配置点内：ODE 残差 =\allowbreak{} 舍入；节点之间：残差 \textasciitilde{} 误差$\cdot$N\textasciicircum{}2(\allowbreak{}插值误差的导数) & N=\allowbreak{}8: 内点残差 1.7e-13，细网格残差 3.8e-01，误差 1.8e-04；N=\allowbreak{}12: 内点残差 1.9e-12，细网格残差 6.1e-04，误差 6.0e-08；N=\allowbreak{}16: 内点残差 2.3e-12，细网格残差 1.0e-07，误差 3.5e-12；N=\allowbreak{}20: 内点残差 5.2e-12，细网格残差 2.9e-11，误差 8.5e-15；细网格残差/\allowbreak{}误差 最小 2104 倍 \checkmark \\
A & 精确解的谱微分残差地板 & 微分放大舍入 \textasciitilde{}N\textasciicircum{}4$\cdot$u(\allowbreak{}谱微分矩阵 ||D\textasciicircum{}2||\textasciitilde{}N\textasciicircum{}4) & N=\allowbreak{}20: 残差 6.82e-12(\allowbreak{}N\textasciicircum{}4u=\allowbreak{}1.8e-11)，插值误差 3.33e-15；N=\allowbreak{}24: 残差 2.10e-11(\allowbreak{}N\textasciicircum{}4u=\allowbreak{}3.7e-11)，插值误差 2.44e-15；N=\allowbreak{}32: 残差 1.08e-10(\allowbreak{}N\textasciicircum{}4u=\allowbreak{}1.2e-10)，插值误差 2.66e-15；实测/\allowbreak{}预测 最大比例 0.93 \checkmark \\
B & 配置法与 tau 法的解之差 & $\le$ 1.00000e-10 & $2.677\times 10^{-11}$ \checkmark \\
A & 配置法 vs tau 法 & 同一多项式空间、\allowbreak{}同一方程：两条路线应当给同一个解 & N=\allowbreak{}16: 配置 3.5e-12 /\allowbreak{} tau 2.7e-11 /\allowbreak{} 互差 2.7e-11；N=\allowbreak{}20: 配置 8.5e-15 /\allowbreak{} tau 6.9e-15 /\allowbreak{} 互差 9.7e-15；N=\allowbreak{}24: 配置 7.5e-15 /\allowbreak{} tau 5.5e-15 /\allowbreak{} 互差 9.1e-15；N=\allowbreak{}32: 配置 1.9e-14 /\allowbreak{} tau 5.4e-15 /\allowbreak{} 互差 2.1e-14 \checkmark \\
B & 谱配置的边界残差 & $\le$ 1.00000e-11 & $1.990\times 10^{-13}$ \checkmark \\
A & 谱配置的边界残差与误差同量级 & 边界行是 D\_N 的一行，解出来时行残差 \textasciitilde{} 舍入(\allowbreak{}与内点同量级) & smooth N=\allowbreak{}16: 离散=\allowbreak{}5.9e-16 真实=\allowbreak{}5.9e-16 误差=\allowbreak{}3.5e-12；smooth N=\allowbreak{}24: 离散=\allowbreak{}4.6e-15 真实=\allowbreak{}4.6e-15 误差=\allowbreak{}7.5e-15；smooth N=\allowbreak{}32: 离散=\allowbreak{}8.8e-15 真实=\allowbreak{}8.8e-15 误差=\allowbreak{}1.9e-14；neumann N=\allowbreak{}16: 离散=\allowbreak{}8.2e-14 真实=\allowbreak{}8.2e-14 误差=\allowbreak{}1.1e-13；neumann N=\allowbreak{}24: 离散=\allowbreak{}8.5e-14 真实=\allowbreak{}8.5e-14 误差=\allowbreak{}6.2e-15；neumann N=\allowbreak{}32: 离散=\allowbreak{}2.0e-13 真实=\allowbreak{}2.0e-13 误差=\allowbreak{}4.2e-13 \checkmark \\
B & 聚集比 / N - 2/pi & $\le$ 0.001 & $3.196\times 10^{-5}$ \checkmark \\
A & 节点聚集的几何量 & 间距比 \textasciitilde{} (\allowbreak{}2/\allowbreak{}pi)N；密度 rho(\allowbreak{}x)=\allowbreak{}N/\allowbreak{}(\allowbreak{}pi sqrt(\allowbreak{}1-x\textasciicircum{}2))，rho(\allowbreak{}0.9)/\allowbreak{}rho(\allowbreak{}0)=\allowbreak{}1/\allowbreak{}sqrt(\allowbreak{}0.19)=\allowbreak{}2.294 & N=\allowbreak{}8: 最大/\allowbreak{}最小间距 =\allowbreak{} 5.027(\allowbreak{}/\allowbreak{}N =\allowbreak{} 0.6284)；N=\allowbreak{}16: 最大/\allowbreak{}最小间距 =\allowbreak{} 10.153(\allowbreak{}/\allowbreak{}N =\allowbreak{} 0.6346)；N=\allowbreak{}32: 最大/\allowbreak{}最小间距 =\allowbreak{} 20.355(\allowbreak{}/\allowbreak{}N =\allowbreak{} 0.6361)；N=\allowbreak{}64: 最大/\allowbreak{}最小间距 =\allowbreak{} 40.735(\allowbreak{}/\allowbreak{}N =\allowbreak{} 0.6365)；N=\allowbreak{}128: 最大/\allowbreak{}最小间距 =\allowbreak{} 81.483(\allowbreak{}/\allowbreak{}N =\allowbreak{} 0.6366)；N=\allowbreak{}64 时 rho(\allowbreak{}0.9)/\allowbreak{}rho(\allowbreak{}0)=\allowbreak{}2.294157(\allowbreak{}解析 2.294157) \checkmark \\
B & Bratu(x$\in$[0,1])的谱配置误差 & $\le$ 1.00000e-11 & $2.220\times 10^{-15}$ \checkmark \\
A & 区间映射的伸缩因子 s=2/(b-a) & d/\allowbreak{}dx =\allowbreak{} s$\cdot$d/\allowbreak{}dxi，d\textasciicircum{}2/\allowbreak{}dx\textasciicircum{}2 =\allowbreak{} s\textasciicircum{}2$\cdot$d\textasciicircum{}2/\allowbreak{}dxi\textasciicircum{}2；[0,\allowbreak{}1] 上 s=\allowbreak{}2 & 正确伸缩：误差 2.22e-15；漏掉伸缩：偏差 7.59e-01(\allowbreak{}s=\allowbreak{}2 时算子整体差 4 倍) \checkmark \\
C & 指数收敛 vs 幂律收敛的判别 & 解析解 err\textasciitilde{}e\textasciicircum{}\{-lambda N\}(\allowbreak{}比值递增)；C\textasciicircum{}s 解 err\textasciitilde{}N\textasciicircum{}\{-s\}(\allowbreak{}比值 $\to$ 常数) & smooth(\allowbreak{}N=\allowbreak{}8..18)：斜率 -2.2568，R\textasciicircum{}2=\allowbreak{}0.99636，误差比 ['39.3',\allowbreak{} '76.0',\allowbreak{} '196.7',\allowbreak{} '88.0']；abs3(\allowbreak{}N=\allowbreak{}8..32)：p=\allowbreak{}1.9729，R\textasciicircum{}2=\allowbreak{}0.999996，误差比 ['1.56',\allowbreak{} '1.43',\allowbreak{} '1.36',\allowbreak{} '1.30',\allowbreak{} '1.26',\allowbreak{} '1.23',\allowbreak{} '1.21',\allowbreak{} '1.19',\allowbreak{} '1.17',\allowbreak{} '1.16',\allowbreak{} '1.15',\allowbreak{} '1.14'] \checkmark \\
C & C\textasciicircum{}s 解的配置率 = 插值率 - 1 & 微分吃掉一阶光滑度：配置率 =\allowbreak{} 插值率 - 1(\allowbreak{}实测 |x|\textasciicircum{}3: 3.00$\to$1.99；|x|\textasciicircum{}5: 5.01$\to$4.02) & abs3: 配置解 p=\allowbreak{}1.987，插值 p=\allowbreak{}3.002，差 1.015(\allowbreak{}理论 1)；abs5: 配置解 p=\allowbreak{}4.023，插值 p=\allowbreak{}5.012，差 0.989(\allowbreak{}理论 1) \checkmark \\
C & 解析算例的舍入地板 & 地板 \textasciitilde{} N\textasciicircum{}4$\cdot$u(\allowbreak{}微分放大舍入)，N$\approx$18 起到达 & 地板 7.550e-15(\allowbreak{}N$\ge$18)，最小误差 24 处 7.55e-15；N=\allowbreak{}64 时 9.26e-14(\allowbreak{}比 N=\allowbreak{}24 的 7.55e-15 更大) \checkmark \\
C & 微分矩阵的条件数增长 & kappa+ \textasciitilde{} N\textasciicircum{}2(\allowbreak{}sigma\_min+ 有界)，朴素 kappa 由零空间(\allowbreak{}常数)决定，无意义 & kappa+ =\allowbreak{} ['35.1',\allowbreak{} '74.8',\allowbreak{} '130.0',\allowbreak{} '286.2',\allowbreak{} '503.9',\allowbreak{} '1123.2',\allowbreak{} '1987.9']；拟合指数 1.946(\allowbreak{}R\textasciicircum{}2=\allowbreak{}0.99989)；朴素 kappa 最大 5.93e+18，比 kappa+ 大 3.0e+15 倍 \checkmark \\
C & 陡峭解析解所需的分辨率 & tanh(\allowbreak{}kx) 在带形内解析 =\allowbreak{}> 仍指数收敛，但速率 \textasciitilde{} e\textasciicircum{}\{-c$\cdot$N/\allowbreak{}k\} & err<1e-10 的最小 N：k=\allowbreak{}3 -> 48，k=\allowbreak{}6 -> 96，k=\allowbreak{}10 -> None(\allowbreak{}窗口上限 128)；半对数斜率 k=\allowbreak{}3: -0.589，k=\allowbreak{}6: -0.276，k=\allowbreak{}10: -0.111 \checkmark \\
\end{longtable}


\section{它会怎样骗你（本章版）}

\begin{enumerate}
  \item \textbf{把灵敏度当充分判据}：同一个 $S=5.885\times10^{15}$，
        增长模态误差 $4.8\times10^{-10}$、衰减模态 $1.7\times10^{1}$。
  \item \textbf{把衰减模态的失效当成"步长不够小"}：$h$ 缩小 32 倍误差不降，
        因为它是舍入放大（\eqref{eq:decay-err}）。
  \item \textbf{把打靶残差小当解准}：$|F|$ 的可分辨下限是 $u\cdot S$，
        $kL=20$ 时 $|F|$ 停在 $8.2\times10^{-10}$ 而解误差同量级。
  \item \textbf{拿"再打一枪"当解析真值}：参考斜率与 Newton 解共享同一套舍入，是循环论证。
  \item \textbf{Jacobian 取步末状态}：二次收敛会悄悄退化成线性收敛
        （残差每步只降 $1.3\times10^{-3}$ 倍），只有把残差史打出来才看得见。
  \item \textbf{用内部阶代表整体阶}：一阶边界格式下内部阶也是 $1.00$，
        误差会铺满整个区间。
  \item \textbf{报"边界残差 $10^{-13}$"}：那是格式自己那一行的残差；
        真实违反量在同一帧里是 $1.9\times10^{-2}$。
  \item \textbf{用线性系统残差当精度}：残差 $8.7\times10^{-14}$ 时误差还有 $1.57\times10^{-4}$。
  \item \textbf{报阶不写窗口}：同一批数据 $1.86$ 与 $1.98$ 都"对"。
  \item \textbf{用绝对残差当 Newton 判据}：残差地板是 $u/h^2$ 量级，绝对判据会让迭代空转。
  \item \textbf{只说"指数收敛"}：没有斜率、没有地板、没有撞地板的 $N$，
        读者无法决定该取多大 $N$；而且地板之后继续加密\textbf{更差}。
  \item \textbf{只看 $R^2$ 判别指数/幂律}：解析解的窗口里幂律也能拟合到 $R^2\approx0.98$；
        要看误差比的走向。
  \item \textbf{拿插值误差冒充配置解误差}：同一批节点上 $3.00$ 与 $1.99$（$|x|^3$），
        差整整一阶。
  \item \textbf{把节点上的 $0$ 当精度}：插值误差在节点上按定义就是 $0$。
  \item \textbf{忘记区间映射}：$[0,1]$ 上漏掉 $s=2$ 会解出偏差 $O(1)$ 的结果；
        识别 Lobatto 网格时也要允许平移缩放，否则边界残差量出 $\pi/2$。
  \item \textbf{把 $D_N$ 直接拿去解方程组}：它稠密、奇异（$D_N\mathbf{1}=0$），
        $\sigma_{\max}/\sigma_{\min}$ 给出 $10^{17}$ 的噪声；
        配置矩阵的条件数按 $N^{4.4}$ 长，$N$ 的上限由它决定。
  \item \textbf{在等距网格上做全局多项式插值}：$n=50$ 时连直线都插不准（误差 $1.5\times10^{-3}$），
        $n=100$ 时是 $3.49$；搬有限差分解要用分段线性。
  \item \textbf{把 tau 法当"另一个配置法"}：它的边界行写法不同，
        在 Neumann 边界上直接套会写错。
\end{enumerate}
